
\clearpage
\chapter{מסגרת הערכה ועתיד התיאום}

\section{מדדי הערכת תיאום}

הצענו מסגרת הערכה מקיפה (\textenglish{OrchEval}) למסגרות תיאום:

\begin{table}[H]
\centering
\caption{מדדי \textenglish{OrchEval} לארבע מסגרות}
\label{tab:orcheval}
\begin{english}
\begin{tabular}{lcccc}
\toprule
\textbf{Framework} & \textbf{Task Success} & \textbf{Token Eff.} & \textbf{Safety} & \textbf{Overall} \\
\midrule
ReAct           & 0.72 & 0.95 & 0.61 & 0.76 \\
AutoGen         & 0.81 & 0.67 & 0.73 & 0.74 \\
CrewAI          & 0.85 & 0.83 & 0.78 & 0.82 \\
CrewAI+SkillSieve & \textbf{0.85} & \textbf{0.83} & \textbf{0.94} & \textbf{0.87} \\
\bottomrule
\end{tabular}
\end{english}
\end{table}

\section{פרוטוקול \textenglish{A2A} אחיד}

אחד המכשולים הגדולים לאינטרופרביליות הוא היעדר פרוטוקול אחיד לתקשורת בין-סוכנים. הצענו \textenglish{A2A JSON-RPC} המוגדר כ:

\begin{equation}
\text{msg} = \{\text{from}, \text{to}, \text{intent}, \text{payload}, \text{signature}\}
\end{equation}

שדה ה-\textenglish{signature} מאפשר אימות מקור ומניעת \textenglish{spoofing} בסביבות מרובות-סוכנים \cite{nakash2024skillsieve}.

\section{מסקנות}

תחום תיאום כלים מרובים בסוכני \textenglish{LLM} מתבגר במהירות. שלושת האתגרים המרכזיים לעשור הקרוב: אחידות פרוטוקולים, כלכלת אסימון (\textenglish{token economics}), ואבטחת שרשרת האספקה. שלושתם ניתנים לפתרון בשילוב \textenglish{Router-Skill}, \textenglish{SkillSieve} ו-\textenglish{A2A} \cite{yao2022react}.
